%=====================================================================
\chapter{Network Virtualization and SDN}\label{ch:network}
%=====================================================================
\locator{4}{Part IV --- Storage, Network and Cloud Virtualization}

\begin{outcomes}
By the end of this chapter you will be able to:
\begin{tight}
  \item \textbf{Describe} a virtual switch and \textbf{explain} what it does
        that a physical switch cannot.
  \item \textbf{State} the limit that VLANs reach, and \textbf{explain} how an
        overlay removes it.
  \item \textbf{Compute} the encapsulation overhead of VXLAN and
        \textbf{predict} the MTU failure it causes.
  \item \textbf{Explain} the control-plane/data-plane split and \textbf{say}
        what software-defined networking actually moved.
  \item \textbf{Diagnose} a network fault in which small packets succeed and
        large transfers hang.
\end{tight}
\end{outcomes}

\begin{prereq}
\chref{paravirt} (virtio-net, SR-IOV, and the fact that VF-to-VF traffic never
reaches the host) and \chref{containers} (the network namespace, which is what
a container's ``own interface'' actually is). From networking you need Ethernet
frames, IP, MTU and what a switch does with a MAC address table.
\end{prereq}

\begin{keyterms}
virtual switch $\bullet$ Linux bridge $\bullet$ Open vSwitch $\bullet$ tap
device $\bullet$ veth pair $\bullet$ uplink $\bullet$ VLAN $\bullet$ 802.1Q
$\bullet$ trunk $\bullet$ overlay $\bullet$ underlay $\bullet$ VXLAN $\bullet$
VNI $\bullet$ VTEP $\bullet$ Geneve $\bullet$ encapsulation $\bullet$ MTU
$\bullet$ MSS clamping $\bullet$ path MTU discovery $\bullet$ control plane
$\bullet$ data plane $\bullet$ SDN $\bullet$ OpenFlow $\bullet$ NFV $\bullet$
VNF $\bullet$ service chaining $\bullet$ CNI $\bullet$ micro-segmentation
\end{keyterms}

\section{The Switch That Is Not There}

A host with thirty guests has one or two physical network ports. Something must
decide which frame belongs to which guest, and that something is a
\term{virtual switch} in software.

Each guest's virtual NIC terminates in a \term{tap device} on the host --- a
virtual interface whose other end is the guest. The switch learns MAC addresses
and forwards between taps and the physical \term{uplink}, exactly as a physical
switch forwards between ports. Linux offers two: the simple \term{bridge}, and
\term{Open vSwitch}, which adds flow tables, tunnelling and a remote control
interface.

\begin{conceptbox}[What a Virtual Switch Can Do That Steel Cannot]
It is not merely a cheaper switch. Three things follow from being software
inside the hypervisor.

\smallskip
\textbf{It knows what each port is.} A physical switch sees a MAC address. A
virtual switch knows that this port is \emph{the web server of tenant 14}, and
can apply a rule written in those terms. This is what makes
\term{micro-segmentation} possible: a firewall rule enforced at every guest's
own interface, instead of at a choke point the traffic may never cross.

\smallskip
\textbf{It moves with the guest.} Live-migrate a guest (\chref{migration}) and
its port, its rules and its counters are recreated on the destination. No
cable is patched and no switch port is reconfigured.

\smallskip
\textbf{It is created and destroyed in milliseconds.} A Kubernetes node
starting a Pod creates a \term{veth pair}, attaches one end to a bridge and
moves the other into the Pod's network namespace. That is the whole of
container networking, and it happens hundreds of times a minute.

\smallskip
The cost: it is a switch made of your host's CPU. At high packet rates it
competes with the guests, which is the measurement \chref{paravirt} made and
the reason SR-IOV exists.
\end{conceptbox}

\section{VLANs, and the Number 4094}

Separating tenants on a shared network is an old problem with an old answer. A
\term{VLAN} tag --- 802.1Q --- is four bytes inserted into the Ethernet header,
and switches treat differently-tagged frames as if they were on different
wires.

The tag contains a 12-bit VLAN identifier. That is $2^{12} = 4096$ values, of
which 0 and 4095 are reserved: \textbf{4094 usable networks}, per layer-2
domain, for ever.

For one company that is ample. For a cloud provider with fifty thousand
customers it is not a limit to be worked around; it is a wall. And there is a
second problem that bites even small estates: a VLAN exists only within one
layer-2 domain, so a guest cannot keep its address when it migrates to a host
in another building.

\section{Overlays}

An overlay solves both by putting the tenant's Ethernet frame \emph{inside} an
ordinary IP packet. The physical network --- the \term{underlay} --- carries
IP between hosts and knows nothing about tenants. The virtual network --- the
\term{overlay} --- lives entirely in the encapsulation.

\begin{definitionbox}[VXLAN, VNI and VTEP]
\term{VXLAN} wraps a guest's complete Ethernet frame in UDP, and sends it to
the host where the destination guest lives.

\smallskip
The \term{VNI} --- VXLAN network identifier --- is \textbf{24 bits}, so
$2^{24} = 16{,}777{,}216$ virtual networks instead of 4094. A \term{VTEP}
(VXLAN tunnel endpoint) is the component that encapsulates and decapsulates;
on a virtualization host it is the virtual switch.

\smallskip
Because the outer packet is ordinary IP, the overlay crosses any routed network
--- another rack, another building, another data centre --- without that
network being told anything. \term{Geneve} is the newer, extensible version and
is what most recent stacks use; the arithmetic below is the same.
\end{definitionbox}

\dsfig{%
\begin{tikzpicture}[font=\scriptsize,
  f/.style={rounded corners=1.5pt, line width=0.7pt, align=center,
            minimum height=0.62cm, inner sep=2pt, font=\scriptsize},
  inner/.style={f, draw=secondary, fill=secondary!10,
                text=secondary!80!black},
  outer/.style={f, draw=purplehl, fill=purplehl!12, text=purplehl!70!black},
  pay/.style={f, draw=neutral, fill=neutral!12, text=neutral!40!black},
  note/.style={font=\scriptsize\itshape, align=center}
]
  % --- the original frame ---
  \node[font=\scriptsize\bfseries, text=secondary!80!black, anchor=west]
       at (0,2.9) {What the guest sent};
  \node[inner, minimum width=2.0cm] at (1.0,2.3) {Eth\\14 B};
  \node[inner, minimum width=2.0cm] at (3.1,2.3) {IP\\20 B};
  \node[inner, minimum width=1.8cm] at (5.1,2.3) {TCP\\20 B};
  \node[pay,   minimum width=5.4cm] at (8.8,2.3) {payload};
  \draw[decorate, decoration={brace, amplitude=3pt}, secondary!70]
        (11.5,2.68) -- (0,2.68);
  \node[note, text=secondary!80!black] at (5.75,3.05) {};

  % --- the encapsulated frame ---
  \node[font=\scriptsize\bfseries, text=purplehl!70!black, anchor=west]
       at (0,1.3) {What goes on the wire};
  \node[outer, minimum width=1.5cm] at (0.75,0.7) {Eth\\14 B};
  \node[outer, minimum width=1.5cm] at (2.3,0.7)  {IP\\20 B};
  \node[outer, minimum width=1.2cm] at (3.65,0.7) {UDP\\8 B};
  \node[outer, minimum width=1.2cm] at (4.9,0.7)  {VXLAN\\8 B};
  \node[inner, minimum width=6.0cm] at (8.55,0.7)
       {the entire original frame, unchanged};

  \draw[decorate, decoration={brace, amplitude=4pt, mirror}, accent,
        line width=0.9pt] (0,0.32) -- (5.5,0.32);
  \node[font=\scriptsize\bfseries, text=accent, anchor=north] at (2.75,0.18)
       {50 bytes of encapsulation};

  \draw[-{Stealth[length=2mm]}, accent!70, line width=0.9pt]
        (5.75,1.95) -- (5.75,1.05);

  % --- the consequence ---
  \node[rounded corners=3pt, draw=accent!60, fill=accent!5,
        text=accent!70!black, font=\scriptsize, align=left, inner sep=5pt,
        text width=11.4cm] at (5.75,-1.0)
       {\textbf{The consequence.} If the underlay MTU is the usual 1500, the
        \emph{inner} frame may be at most $1500 - 50 = 1450$ bytes. A guest
        still configured for 1500 will send frames that do not fit --- and
        because the outer packet normally has the \emph{don't fragment} bit
        set, they are dropped rather than split.};
\end{tikzpicture}}%
{VXLAN encapsulation, to scale. The 24-bit VNI in those eight bytes is what
replaces the 12-bit VLAN tag; the 50 bytes are what causes the fault in the
worked example.}{vxlan}

\begin{worked}[Small Packets Work, Large Transfers Hang]
A tenant reports: SSH logs in and then freezes on the first command with any
output. \texttt{ping} succeeds. Database queries returning one row work;
queries returning a thousand rows time out. Backups never complete. The network
team reports no errors and no packet loss.

\smallskip
\textbf{Step 1 --- what the two working cases have in common.} \texttt{ping}
sends 64-byte packets. An SSH login exchanges a few hundred bytes. A one-row
query is small. Everything that \emph{works} is small, and everything that
fails is large. That is an MTU fault, and it is almost the only thing that
produces this exact pattern.

\smallskip
\textbf{Step 2 --- the arithmetic.} The overlay is VXLAN over a standard
network:
\[ \underbrace{14}_{\text{outer Eth}} + \underbrace{20}_{\text{outer IP}}
   + \underbrace{8}_{\text{UDP}} + \underbrace{8}_{\text{VXLAN}}
   = 50 \text{ bytes} \]
\[ \text{largest inner frame} = 1500 - 50 = \mathbf{1450 \text{ bytes}} \]

\smallskip
\textbf{Step 3 --- what the guest believes.} Its interface is still at the
default MTU of 1500. TCP therefore negotiates a maximum segment size of
\[ 1500 - \underbrace{20}_{\text{IP}} - \underbrace{20}_{\text{TCP}}
   = 1460 \text{ bytes} \]
so a full segment becomes a 1500-byte inner frame, which becomes a 1550-byte
packet on a wire that carries 1500.

\smallskip
\textbf{Step 4 --- why it hangs rather than fails.} The VTEP sets the
\emph{don't fragment} bit, so the oversized packet is dropped and an ICMP
``fragmentation needed'' message is returned to tell the sender to use a smaller
size. That mechanism --- path MTU discovery --- is how this is supposed to
self-correct.

\smallskip
It does not self-correct, because ICMP is very commonly filtered. The sender
never learns, retransmits the same oversized segment, and the connection stalls
for ever. This is a \term{PMTUD black hole}, and the handshake succeeded only
because the handshake packets were small.

\smallskip
\textbf{Step 5 --- three fixes, in order of preference.}
\begin{tight}
  \item \textbf{Raise the underlay.} Set the physical network to jumbo frames,
        9000 bytes, giving an inner MTU of 8950. Guests keep 1500 and nothing
        else changes. Best answer, and it requires every switch on the path.
  \item \textbf{Lower the guest.} Set guest interfaces to 1450. Correct, free,
        and requires touching every guest --- and every guest anyone adds
        later.
  \item \textbf{Clamp the MSS.} Have the virtual switch rewrite the MSS in
        passing SYN packets to 1410. Works without touching guests, and only
        helps TCP: UDP and anything tunnelled stays broken.
\end{tight}

\smallskip
\textbf{Step 6 --- what to take from it.} The encapsulation of this chapter is
invisible until it is not, and when it becomes visible the symptom points
nowhere near the cause. ``Small works, large hangs'' should go straight to
arithmetic: add up the headers.
\end{worked}

\section{Separating the Planes}

Every switch does two different jobs. The \term{data plane} forwards frames, at
line rate, millions per second. The \term{control plane} decides \emph{how}
they should be forwarded --- learning addresses, running routing protocols,
holding the policy.

Traditionally both live in every switch, so policy is configured box by box and
the network's behaviour is an emergent property of several hundred
configurations nobody can read at once.

\begin{definitionbox}[Software-Defined Networking]
\term{SDN} separates the two: the data plane stays in the switches, and the
control plane is lifted out into a \term{controller} with a view of the whole
network. It programmes each device's forwarding tables through an open
interface --- \term{OpenFlow} being the first and best known.

\smallskip
The change is in where policy lives. ``Tenant 14 may not reach tenant 9'' stops
being a set of access lists on twelve switches and becomes one statement to one
controller, which works out the consequences everywhere. Networks become
programmable in the way that \chref{management}'s infrastructure-as-code makes
servers programmable.
\end{definitionbox}

\begin{conceptbox}[What Actually Happened to SDN]
It is worth being honest about this, because the textbook account and the
industry's experience differ.

\smallskip
Pure OpenFlow --- dumb switches, one central controller --- did not take over.
The controller became a scaling and availability problem, hardware support was
uneven, and operators were unwilling to put every forwarding decision behind a
single piece of software.

\smallskip
What did take over is the \emph{idea}: a central place that holds intent, and
devices programmed from it. It is everywhere in virtualization --- NSX, ACI,
every cloud's virtual network, every Kubernetes CNI plugin --- usually with
distributed agents rather than one controller, and usually with the overlay
doing the heavy lifting. The vocabulary survived better than the protocol,
which is a common pattern and worth recognising when you meet the next one.
\end{conceptbox}

\section{NFV and Container Networking}

If a switch can be software, so can a firewall, a load balancer and a router.
\term{Network function virtualization} replaces purpose-built appliances with
virtual machines or containers --- \term{VNFs} --- running the same function.
A chain of them, traffic steered through in order, is \term{service chaining},
and the economics are the ones from \chref{what}: appliances sized for peak and
idle most of the time, consolidated.

In containers the same job is done by \term{CNI}, the Container Network
Interface. When a Pod starts, the kubelet calls a CNI plugin, which creates the
veth pair, assigns an address and installs routes. Which plugin you choose
decides whether you get an overlay (Flannel's VXLAN), routed addresses without
encapsulation (Calico in BGP mode), or an eBPF data path with policy (Cilium).

\begin{alertbox}[Where the packets you cannot see go]
Three places in this book where traffic exists and your monitoring does not see
it, each with a different cause:

\smallskip
\textbf{Guest to guest on one host.} The frame goes guest $\rightarrow$ virtual
switch $\rightarrow$ guest and never touches the uplink. A physical tap or a
switch SPAN port sees nothing.

\smallskip
\textbf{VF to VF on one SR-IOV card} (\chref{paravirt}). The card's internal
switch forwards it, so even the hypervisor is not involved.

\smallskip
\textbf{Pod to Pod on one node.} The veth pair and the node's bridge handle it
entirely.

\smallskip
In each case the fix is the same in principle --- the observation point must be
the virtual switch, not the wire --- and it is the reason
\chref{security}'s detection chapter cares so much about where sensors are
placed.
\end{alertbox}

\begin{pitfall}
\begin{tight}
  \item \textbf{Deploying an overlay without fixing MTU.} The worked example.
        Decide between jumbo underlay, lowered guest MTU and MSS clamping
        \emph{before} the first tenant, not after the first backup fails.
  \item \textbf{Filtering ICMP wholesale.} It breaks path MTU discovery and
        turns a clean error into a silent hang. Permit type 3 code 4.
  \item \textbf{Assuming the overlay is free.} Encapsulation costs CPU on every
        packet in both directions. Check whether your NIC offloads VXLAN; the
        difference is large.
  \item \textbf{Expecting physical network monitoring to see virtual
        traffic.} Most east--west traffic never reaches a physical port.
  \item \textbf{Treating 4094 as a comfortable number.} It is per layer-2
        domain and it also constrains the \emph{underlay} design, not just
        tenant count.
  \item \textbf{Choosing a CNI plugin by popularity.} Overlay, routed and eBPF
        plugins differ in MTU behaviour, in policy support and in how they fail.
        That choice is harder to change than any other in a cluster.
\end{tight}
\end{pitfall}

\begin{lab}[Building an Overlay, and Breaking It]
Builds a VXLAN overlay between two network namespaces, then reproduces the
black hole of the worked example and fixes it three ways. Needs one Linux host.
Expect forty minutes.

\begin{lstlisting}[language=bash]
#!/usr/bin/env bash
# ch13_overlay.sh -- a virtual switch, a VXLAN tunnel and an MTU black hole.
set -euo pipefail

# --- 1. A virtual switch, and two "guests" as namespaces. --------------
sudo ip link add br-vss type bridge
sudo ip link set br-vss up
for n in 1 2; do
  sudo ip netns add "g$n"
  sudo ip link add "veth$n" type veth peer name "in$n"
  sudo ip link set "veth$n" master br-vss
  sudo ip link set "veth$n" up
  sudo ip link set "in$n" netns "g$n"
  sudo ip netns exec "g$n" ip addr add "10.10.0.$n/24" dev "in$n"
  sudo ip netns exec "g$n" ip link set "in$n" up
  sudo ip netns exec "g$n" ip link set lo up
done
# This is container networking, in nine lines.
sudo ip netns exec g1 ping -c 2 10.10.0.2
bridge fdb show br br-vss | head -4     # the switch has learned MACs

# --- 2. Now the underlay: two more namespaces playing "hosts". --------
sudo ip netns add h1 ; sudo ip netns add h2
sudo ip link add u1 type veth peer name u2
sudo ip link set u1 netns h1 ; sudo ip link set u2 netns h2
sudo ip netns exec h1 ip addr add 192.168.50.1/24 dev u1
sudo ip netns exec h2 ip addr add 192.168.50.2/24 dev u2
for h in h1 h2; do
  sudo ip netns exec "$h" ip link set lo up
done
sudo ip netns exec h1 ip link set u1 up
sudo ip netns exec h2 ip link set u2 up
sudo ip netns exec h1 ping -c 2 192.168.50.2

# --- 3. A VXLAN tunnel across it: VNI 100. ----------------------------
sudo ip netns exec h1 ip link add vx0 type vxlan id 100 \
     remote 192.168.50.2 local 192.168.50.1 dstport 4789 dev u1
sudo ip netns exec h2 ip link add vx0 type vxlan id 100 \
     remote 192.168.50.1 local 192.168.50.2 dstport 4789 dev u2
sudo ip netns exec h1 ip addr add 172.16.0.1/24 dev vx0
sudo ip netns exec h2 ip addr add 172.16.0.2/24 dev vx0
sudo ip netns exec h1 ip link set vx0 up
sudo ip netns exec h2 ip link set vx0 up
sudo ip netns exec h1 ip -d link show vx0 | grep -o 'vxlan id 100'

# --- 4. Small packets work. ------------------------------------------
sudo ip netns exec h1 ping -c 2 172.16.0.2

# --- 5. Now the black hole: large packets, don't-fragment set. --------
# The underlay is 1500; VXLAN needs 50; so 1450 is the real limit.
sudo ip netns exec h1 ip link set vx0 mtu 1500    # deliberately wrong
echo "--- 1400 bytes: fits ---"
sudo ip netns exec h1 ping -c 2 -M do -s 1400 172.16.0.2 || true
echo "--- 1472 bytes: does not fit, and will not say so ---"
sudo ip netns exec h1 ping -c 2 -W 2 -M do -s 1472 172.16.0.2 || \
  echo "no reply -- exactly the worked example's symptom"

# --- 6. Fix 1: lower the overlay MTU to the arithmetic. ---------------
sudo ip netns exec h1 ip link set vx0 mtu 1450
sudo ip netns exec h2 ip link set vx0 mtu 1450
sudo ip netns exec h1 ping -c 2 -M do -s 1422 172.16.0.2

# --- 7. Fix 2: raise the underlay instead, and keep 1500 inside. ------
sudo ip netns exec h1 ip link set u1 mtu 9000
sudo ip netns exec h2 ip link set u2 mtu 9000
sudo ip netns exec h1 ip link set vx0 mtu 8950
sudo ip netns exec h2 ip link set vx0 mtu 8950
sudo ip netns exec h1 ping -c 2 -M do -s 8000 172.16.0.2

# --- 8. Watch the encapsulation on the wire. -------------------------
sudo timeout 6 ip netns exec h2 tcpdump -ni u2 -c 4 udp port 4789 &
sleep 1
sudo ip netns exec h1 ping -c 3 172.16.0.2 >/dev/null ; wait || true
# Each capture line is one outer packet carrying one whole inner frame.

# --- 9. Tear down. ---------------------------------------------------
for n in g1 g2 h1 h2; do sudo ip netns del "$n" 2>/dev/null || true; done
sudo ip link del br-vss 2>/dev/null || true
\end{lstlisting}

\textbf{What to notice.} Step 5 is the fault that costs the most time in
production, reproduced in two commands, and the difference between the working
and failing case is 72 bytes. Step 8 shows why: the capture is full-size outer
packets, each carrying a frame that was never meant to be carried.
\end{lab}

\begin{checkpoint}
\begin{tightnum}
  \item An overlay uses Geneve with 16 bytes of options, over IPv6. Outer
        headers are Ethernet 14, IPv6 40, UDP 8, Geneve $8 + 16$. Give the
        total overhead and the largest inner frame on a 1500-byte underlay.
  \item Why does a PMTUD black hole present as a hang rather than an error,
        and what single firewall change makes it present honestly?
  \item A tenant must be isolated from all others, and their guests are spread
        over three buildings on different IP subnets. Explain why VLANs cannot
        do this and what you would use.
\end{tightnum}
\end{checkpoint}

\begin{chaptersummary}
\begin{tight}
  \item A virtual switch forwards between guests in software. It can express
        rules about identity rather than MAC addresses, it moves with the
        guest, and it is made of your host's CPU.
  \item A VLAN tag is 12 bits: 4094 networks, confined to one layer-2 domain.
  \item VXLAN puts the guest's whole frame inside UDP. The VNI is 24 bits ---
        16.7 million networks --- and the outer packet is ordinary IP, so the
        overlay crosses any routed network.
  \item Encapsulation costs 50 bytes, so the inner MTU is 1450 on a 1500-byte
        underlay. Leave the guest at 1500 and large transfers hang silently.
  \item Fix it by raising the underlay to jumbo frames, lowering guest MTU, or
        clamping MSS --- in that order of preference.
  \item SDN separates the control plane from the data plane. Pure OpenFlow did
        not win; the idea of central intent with programmed devices did.
  \item Most east--west traffic never reaches a physical port, so monitoring
        must sit at the virtual switch.
\end{tight}
\end{chaptersummary}

\begin{reviewq}
  \item \textbf{(MCQ)} The VXLAN network identifier is: (a)~12 bits, giving
        4094 networks; (b)~16 bits; (c)~24 bits, giving about 16.7 million;
        (d)~32 bits.
  \item \textbf{(MCQ)} VXLAN encapsulation on an Ethernet/IPv4 underlay adds:
        (a)~4 bytes; (b)~24 bytes; (c)~50 bytes; (d)~100 bytes.
  \item \textbf{(Short)} Give two things a virtual switch can do that a
        physical switch cannot, with the reason each is possible.
  \item \textbf{(Short)} Explain why a VLAN cannot stretch a tenant's network
        across two routed sites, and how an overlay can.
  \item \textbf{(Short)} State the three MTU fixes in order of preference, with
        the drawback of each.
  \item \textbf{(Applied)} A cluster is being built on a network whose switches
        cannot do jumbo frames, with guests managed by a team that will not
        change their images. Choose an approach, state exactly what you would
        configure and where, and name the traffic your choice will not
        protect.
  \item \textbf{(Applied)} A security team asks why their intrusion sensor on
        the top-of-rack switch sees almost no traffic between application
        servers that are known to be busy. Explain the three mechanisms that
        could be responsible and say how you would confirm which it is.
\end{reviewq}
